\documentclass[11pt]{article}
\usepackage[left=1in, right=1in, top=1in, bottom=1in]{geometry}
\usepackage{times}
\usepackage{amsmath}
\usepackage{hyperref}
\usepackage{enumitem}

\begin{document}
\noindent Dear Area Chair and Reviewers,\\

We have carefully considered each comment and made substantial revisions to address the concerns raised. Below, we provide an overview of the main changes, followed by a point-by-point response to each reviewer.\\

\noindent \textbf{Overview of Changes} (corresponding to the Area Chair's suggestions):
\begin{enumerate}[label=(\arabic*), itemsep=1pt]
    \item Modern backbone (Qwen3-4B) added to the main results, best or tied-best on all 9 benchmarks. [Table 1, Section 5.2]
    \item Quantitative compute-efficiency analysis: per-step token and wall-clock breakdown on four model/task settings. [Appendix D]
    \item Direct pivot-localization-accuracy evaluation against oracle labels across three agentic benchmarks. [Appendix H.2]
    \item Evaluation protocol clarification that resolves the apparent baseline-discrepancy concern via a format-vs-distribution-shift decomposition. [Section 5.1, Appendix J]
    \item Additional ablations: retry count, group size, and a complexity-drop study isolating Positive Amplification. [Appendix C]
    \item Structured four-axis comparison with related methods (Critique-GRPO, Reflect-GRPO, PRM, GiGPO, VinePPO, BAPO, GSPO, OPMD). [Appendix E]
\end{enumerate}

\noindent \textbf{Point-by-Point Response:}\\

\noindent \textbf{Reviewer iAPy:}
\begin{enumerate}[label=(\arabic*), itemsep=2pt]
    \item \textbf{[Novelty.]} Appendix E consolidates the design distinctions along four axes (exploration mechanism, credit granularity, external supervision, rollout cost, hyperparameter count). R$^3$L is the only method in the comparison that simultaneously achieves step-level credit without external supervision, keeps rollouts within $N$, and introduces a single scalar $\alpha$ (while dropping $\beta$ and $\epsilon$ from GRPO). Table 2 confirms the three components provide independent and additive gains. [Appendix E, Table 2]

    \item \textbf{[No direct pivot-accuracy evaluation.]} Added in Appendix H.2. ScienceWorld provides automatic oracle labels via step-level subtask rewards; ALFWorld and WebShop oracles are annotated on 100 failed trajectories per benchmark via human judgment plus DeepSeek labeling. Oracle Agreement rises from 36--57\% at step 100 to 61--76\% at step 400, and retry success rate conditioned on a correct pivot is consistently 2--3$\times$ higher than on a wrong pivot (e.g., 85\% vs 43\% on ALFWorld 7B at step 400). [Appendix H.2]

    \item \textbf{[No wall-clock or token-budget analysis.]} Added in Appendix D as a structural rollout-cost table plus per-training-step measurements on four settings. R$^3$L's rollout time is lowest across all four; training tokens on multi-step tasks are roughly halved (e.g., 22{,}300 vs GRPO 44{,}141 on Qwen2.5-1.5B/ALFWorld); Critique-GRPO is 1.8--2.0$\times$ slower due to its $2N$ rollouts. The Limitations section now cites these figures explicitly. [Appendix D, Limitations]
\end{enumerate}

\noindent \textbf{Reviewer 4viW:}
\begin{enumerate}[label=(\arabic*), itemsep=2pt]
    \item \textbf{[System complexity and rollout overhead.]} The structural comparison shows R$^3$L's rollout budget stays within $N$ (vs $2N$ for Critique-GRPO), and per-step measurements confirm rollout time is lower than GRPO's across all four tested settings. Removing KL additionally eliminates the reference-model forward pass. [Appendix D]

    \item \textbf{[Cold-start for small models.]} For Qwen2.5-1.5B the cold-start phase lasts approximately 80 steps, under 20\% of the 400-step training budget. During this phase R$^3$L does not degrade: retries driven by inaccurate reflection receive negative advantages and are suppressed, while group-level contrastive learning continues on the base trajectories. [Limitations, Section 5.4]

    \item \textbf{[Reflection errors reinforcing wrong behavior.]} Three safeguards layered: (a) auxiliary SFT ingests only reflection-retry pairs where the retry strictly improves over the base, filtering misdiagnoses at the data level; (b) distilled retry trajectories with negative advantages are suppressed by group-relative normalization; (c) Positive Amplification weights only successful retries. Empirically, roughly 35\% of retries on 1.5B fail to improve, yet robust gains are obtained across all benchmarks. [Appendix H]

    \item \textbf{[Explicit compute comparison.]} We address this concern in detail in our response to Reviewer iAPy's weakness (3); in short, Appendix D provides a structural rollout-cost table together with per-training-step wall-clock and token measurements on four model/task settings. [Appendix D]

    \item \textbf{[Pivot quality analysis.]} We address this concern in detail in our response to Reviewer iAPy's weakness (2); in short, Appendix H.2 reports oracle-agreement trends over training and retry-success rates conditioned on whether the pivot matches the oracle, across all three agentic benchmarks. [Appendix H.2]

    \item \textbf{[Sensitivity for $\alpha$, group size, retry count.]} $\alpha = 3.0$ is robust across all models and tasks; pushing $\alpha > 5.0$ causes overfitting (Appendix C.1). Under $N \in \{4, 8, 16\}$, R$^3$L leads every $(N, \text{benchmark})$ setting on five benchmarks (new Appendix C.4). One retry per base is optimal: two or four sequential retries hurt (0.928 $\to$ 0.913 $\to$ 0.896 on ALFWorld) because they dilute base diversity and because each reflection is an independent one-shot diagnosis (new Appendix C.3). [Appendix C]
\end{enumerate}

\noindent \textbf{Reviewer XSY9:}
\begin{enumerate}[label=(\arabic*), itemsep=2pt]
    \item \textbf{[Baseline numbers lower than Qwen2.5 technical report.]} Our numbers are \emph{post-RL-training} under the \texttt{<think>/<answer>} format (following DeepSeek-R1 and DAPO), whereas the technical report reports \emph{zero-shot} under \texttt{\textbackslash boxed\{\}}; the two differ along two independent dimensions. For Qwen2.5-1.5B GSM8K specifically, the 0.258 gap between our GRPO (0.474) and the report (0.732) decomposes into a 0.073 format effect (about 28\%) and a 0.191 cross-domain-training effect (about 72\%). In-domain training controls (GSM8K-on-GSM8K; MATH-on-Math500) show all methods improve over zero-shot under both formats, confirming the degradation is distribution shift, not methodological. Section 5.1 now states the protocol explicitly, and Appendix J provides the full decomposition with four-backbone zero-shot tables under both formats. [Section 5.1, Appendix J]

    \item \textbf{[Legacy backbones (Qwen2.5, Llama-3.2).]} Qwen3-4B added to Table 1 as a fourth backbone. R$^3$L attains best or tied-best on all 9 benchmarks. On this stronger backbone, GRPO already reaches 0.934 on GSM8K and GSPO reaches 0.722 on Math500, yet R$^3$L still lifts these to 0.948 and 0.753 respectively, indicating the structural bottlenecks persist across model generations rather than being resolved by scale. [Table 1, Section 5.2]

    \item \textbf{[Multiple moving parts; can Positive Amplification alone match full R$^3$L?]} The complexity-drop study in Appendix C.1 shows GRPO+PA alone reaches 0.807 on ALFWorld (vs GRPO 0.720, R$^3$L 0.928): Positive Amplification contributes roughly 40\% of R$^3$L's total gain over GRPO, while Reflect-then-Retry combined with Pivotal Credit Assignment provides the remaining 60\%. R$^3$L also drops $\beta$ (KL) and $\epsilon$ (IS clipping) from GRPO, so the net change in hyperparameter count is small. [Appendix C.1]

    \item \textbf{[Compute efficiency under equal budget.]} Appendix D reports per-step token counts and wall-clock on four settings. R$^3$L's rollout time is consistently below GRPO's; training tokens on multi-step tasks are roughly halved; Critique-GRPO is 1.8--2.0$\times$ slower due to its $2N$ rollouts. The KL-free objective further eliminates the reference-model forward pass, which is why the rollout-time advantage holds even on single-step math tasks where rollout-turn savings are minimal. [Appendix D]
\end{enumerate}

\vspace{0.1in}
\noindent Thank you once again for your time and consideration.

\vspace{0.1in}
\noindent Sincerely,\\
The Authors

\end{document}
